\section{Initialization of the Brownian regulator}
\label{sec:initialization}

The exact step of Theorem~\ref{thm:rg-step} compares two densities once both
belong to the retained class.  The Wilson density has an elementary
initial representation.  The Brownian density requires one preliminary
integration: we retain the spins at integer horizontal coordinates and
integrate the paths between consecutive columns.  This section proves
that this integration produces the same class, with the strip kernel
\eqref{eq:R3} as its free quadratic part.  We keep the scalar and the finite
Taylor coefficients, since both will be needed after the comparison.

Throughout this section the scale is the initial scale $j=N$.  To avoid
confusing the RG depth with a norm of a Brownian path, we denote the latter
by $\mathcal N$.  Write
\[
 g=b^{-1/2},\qquad p=p_N,\qquad t_*=t_N,\qquad M=M_N,\qquad R=R_N.
\]
Thus $t_*\asymp gp$ and $R\asymp1$.  All constants are uniform when
$b=\beta+O_L(1)$ and $r=1+O_L(b^{-1})$; the local construction also works
uniformly for $r$ in a fixed sufficiently small neighborhood of $1$.
The norms, complete supports, inventories, numerical stencils and
normalization conventions are those of
Definition~\ref{def:rg-retained-class}.

\begin{proposition}[Initial densities]
\label{prop:init-density}
For the parameter choices of Theorem~\ref{thm:rg-step}, the flat Wilson
reference and the Brownian strip integral have exact representations in
the retained class at scale $N$, on every admitted rectangle and with
either time seam.  The Wilson reference starts with zero canonical and
regular-error lists and with flat covering labels.  The Brownian
representation starts from the uncorrected affine couplings
\[
 (B_x^q,B_t^q)=b(r^{-1},r),\qquad
 (B_x^w,B_t^w)=bR^{-2}(r^{-1},r).
\]
Its actual affine couplings differ from these by at most $C_L$.  Their
corrections, and its initial canonical coefficients, are given by the
finite Taylor calculation through normalized order four, with coupling
error $O_L(H_N^{-1.05})$.  Its canonical caps can be fixed before choosing
$P_0$ and $H$.  All assertions include the normalized derivatives and
period prescriptions of the retained class.  The scalar extracted per
initial site is the quantity in \eqref{eq:R14} below.
\end{proposition}

\subsection{The reference density and the exact path decomposition}

For the reference use strictly nearest-neighbor numerical rows at depth
zero.  Enforce $dk=0$ on every plaquette.  A nonflat plaquette has an edge
with $r_e\ge P_N/4$, so this restriction can be read by a covering label
owning a flag.  At a flag, the original quadratic energy exceeds the
clipped direct extraction by at least $cH_Nr_e^2$; every other deleted
row is a nonnegative square.  Group flags into the connected primary
footprints of Section~\ref{sec:rg}, retaining every mask query.  Mask
centers reading a flag can be included through the numerical radius
$c_LM_N$; this costs bounded load per flag.  The reserve
$cH_Nt_*^2\asymp p^2$ pays for the support and positioning prices.
The resulting covering weights satisfy their cap and vanish whenever a
plaquette containing an inventory edge is nonflat.  There are no
canonical or regular errors.  This proves the reference assertion of
Proposition~\ref{prop:init-density}.

For the Brownian representation consider the strip between columns $x$
and $x+1$.  Conditional heat paths on its rows are taken on the real
horizontal lift.  Thus each horizontal integer specifies the true $v$
endpoint and the corresponding sign of the quaternion endpoint.  At a
vertical endpoint bond choose the integer for which $|u|\le P_N/2$,
with symmetric treatment of ties.  This is an exact representation of
the quotient spin, not a restriction of its measure.  We use unnormalized
conditional heat-path measures, disintegrated as densities relative to
$d\omega_3(q)\,dw$.  The reciprocal changes in Haar and heat-kernel
normalizations agree with the loop convention of
Section~\ref{sec:brownian}.

Recall the conversion from the Brownian spatial coordinate to strip units.
One horizontal cell has heat length $r/b$, the fugacity is $b^2/2$, and
$\operatorname{Tr}(G G^{\prime *})=2e^{iU}q\cdot q'$ with the chosen
orientation.  Their product is $br\,ds$.  Extract $br$ per site.
On a vertical bond the remaining interaction is
\begin{equation}
 br\int_0^1\bigl(e^{iU(s)}q(s)\cdot q'(s)-1\bigr)\,ds,
 \label{eq:init-vertical-potential}
\end{equation}
where $U$ is the difference of the two real $v$ paths after aligning their
starting sheets by the vertical integer; the quaternion paths use the
same sign and the fixed seam transport.  Either sign of the phase gives
the same construction.  Also subtract the imaginary linear term
$ibr\int_0^1U(s)\,ds$ and retain it as the global phase.

Here the choice of starting sheets can be made invariant under horizontal
reflection.  Average the starting-column convention and the ending-column
convention, transporting sheets horizontally in the latter.  The two
conventions agree on every strip cell on which the local charts are used.
For their global sums, path fluctuations telescope around each vertical
circle.  The starting convention gives $\pi$ times the sum of vertical
integers at column $x$, and the ending convention gives the analogous
sum at $x+1$.  On summing over strips, their average therefore gives
exactly
\[
 i(br/R)\sum_{e_t}u_{e_t}.
\]
This also holds for nonflat integer configurations.  It both identifies
the initial phase and preserves the reflection convention of the class.

\subsection{Brownian localization and endpoint derivatives}

We first give the path estimates used in the strip integration.  They
serve two purposes: they justify Taylor expansion on small paths and
retain an integrable reserve on its complement.

Let $X$ be four-dimensional Brownian motion on $[0,1]$.  Its first three
components develop into a quaternion path by the Stratonovich equation
\[
 dq=\sqrt r\,g\sum_{a=1}^3 e_aq\circ dX^a,
 \qquad dv=\sqrt r\,g\,dX^4,
\]
where $e_1,e_2,e_3$ are the unit imaginary quaternions.  The starting point
is prescribed.  We use the geometric Stratonovich lift
$(X,\mathbb X)$, with $\mathbb X_I$ the second iterated integral on an
interval $I$.  Fix $1/3<\nu<1/2$ and a sufficiently large integer $m_0$.
If $\mathcal D_n$ is the family of dyadic intervals of length $2^{-n}$,
put
\begin{equation}
 \mathcal N(X)^{4m_0}
 =\sum_{n\ge0}\ \sum_{I\in\mathcal D_n}
 2^{4m_0\nu n}
 \left(|\Delta_I X|^{4m_0}+|\mathbb X_I|^{2m_0}\right).
 \label{eq:init-rough-size}
\end{equation}
The second level is thus measured with its square root.  Smooth profiles
will be functions of the polynomial quantity $\mathcal N^{4m_0}$.

\begin{lemma}[Path and endpoint estimates]
\label{lem:init-path-estimates}
The quantity \eqref{eq:init-rough-size} controls a $\nu$-H\"older rough-path
norm, and
\[
 \Pr\{\mathcal N>D\}\le C e^{-cD^2}.
\]
The same bound holds for a Gaussian bridge.  Adding a path with bounded
first derivative, or applying a smooth deterministic orthogonal rotation
to the driving increments, preserves these estimates with the corresponding
deterministic norm costs.  Every fixed derivative of a smooth profile
under these transformations has a polynomial bound in those norms and
$\mathcal N$.

For fixed endpoint sheets, let $K$ denote the unnormalized conditional
row measure.  Its total density obeys
\begin{align}
 K(1)&\le Cg^{-4}\exp(-cg^{-2}r_{e_x}^2),
 \label{eq:init-heat-bound}\\
 K(1_{\{\mathcal N>D\}})&\le Cg^{-4}e^{-cD^2},
 \label{eq:init-conditioned-tail}
\end{align}
uniformly in the endpoints.  The versions of these bounds needed for
smooth profile and path factors admit the fixed endpoint derivatives of
the retained norm.  Derivatives cost fixed polynomials in
$g^{-1},\mathcal N,M,p$ and leave a positive fraction of the exponential
reserves.
\end{lemma}

\begin{proof}
Choose $\nu<\nu'<1/2$ with $4m_0(\nu'-\nu)>1$.  First-level increments
have Gaussian tails at scale $2^{-n/2}$.  Second-level increments are
second Gaussian chaoses, including their deterministic symmetric part;
their square roots have Gaussian tails at the same scale.  At the
$\nu'$-H\"older normalization the tail at level $n$ is bounded by a
constant times
\[
 \exp\{-cD^2 2^{(1-2\nu')n}\}.
\]
Summing over the $2^n$ intervals first and then over levels proves a
Gaussian tail for the corresponding dyadic supremum.  The condition on
$m_0$ sums the extra weights in \eqref{eq:init-rough-size}.  Dyadic
concatenation then controls arbitrary increments and gives the asserted
rough-path bound.  Subtracting the linear endpoint interpolation gives
the bridge estimate.

For a deterministic path $h$, the first-level increment changes by at
most $\|h'\|_\infty|I|$.  The new second-level terms are bounded by this
quantity times the oscillation of $X$ on $I$, together with a term
$C\|h'\|_\infty^2|I|^2$.  These estimates are summable with the same
weights.  Under $dX\mapsto O(s)dX+h'(s)ds$, subtract on each interval the
frozen value of $O$ times its first increment and the corresponding
tensor square times its second increment.  Integration by parts with
the smooth coefficient gives an additional factor $|I|$ in the
remainders.  This also proves convergence of the differentiated dyadic
tests.  The same estimates apply at every fixed derivative order.

The extension theorem for geometric rough paths
\cite[Theorem~2.2.1]{Lyons1998} bounds higher iterated integrals by the
usual factorial estimates.  Applied to the linear
quaternion equation, it gives a convergent series in the path scale.
The series is also analytic in added constant-velocity parameters, by
augmenting the path with time.  On a set where $g\mathcal N$ is sufficiently
small, the quaternion and circle displacements are bounded by
$Cg\mathcal N$.  Outside that set compactness gives the needed bound for
the quaternion displacement; the circle displacement remains linear.
These statements can equivalently be obtained by dyadic polygonal
Brownian approximations.  Their Stratonovich lifts converge almost surely
in H\"older norm and have a common rough-norm bound with a Gaussian
exponential moment
\cite[Propositions~19 and~26, Corollary~27]{FrizVictoir2005}.
Continuity of extension \cite[Theorem~2.2.2]{Lyons1998}, together with
the factorial bounds, then gives convergence of every fixed iterated
integral in every finite $L^q$.

The heat-kernel Gaussian bound on the fixed compact sphere, multiplied
by the real Gaussian density in the circle lift, gives
\eqref{eq:init-heat-bound}; chord and geodesic distances are comparable
in the required direction.  For \eqref{eq:init-conditioned-tail}, split
the row at half-time.  A large whole-path rough size forces a large size
on one half, up to a fixed factor, by concatenation.  For the first half,
use its Gaussian tail and the heat-kernel supremum on the other half.
For the second half use heat-bridge reversal.  Stratonovich
anti-development reverses with the path, so the same estimate applies.

We spell out the endpoint differentiation, since the unnormalized
measure is important here.  To vary its endpoints, multiply the
quaternion path on the left by a smooth deterministic path $H(s)$ whose
endpoint values make the prescribed changes.  Translate the circle
path in the same way.  The driving increments become exactly
\[
 dX^q\longmapsto \operatorname{Ad}_{H(s)}dX^q
       +\frac{H'(s)H(s)^{-1}}{\sqrt r\,g}\,ds.
\]
The rotation is deterministic and orthogonal.  Its composition with
the deterministic drift changes Brownian measure by the
Cameron--Martin exponential
\cite[Theorem~2]{CameronMartin1944}.  Differentiating that exponential gives
scores bounded by fixed polynomials in $\mathcal N$ and $g^{-1}$.
The preceding transformation estimates control the differentiated path
factors and profiles.  Haar and Lebesgue endpoint translations have
Jacobian one.  H\"older's inequality, using a sufficiently small part of
the Gaussian reserve for the scores, leaves the asserted tail bounds.
Off the path-size profiles, the deleted-row factors use the paths only
through $H(s)q(s)$ and deterministic circle translation, whose derivatives
are bounded; no small-path expansion is used there.

These identities can first be read as identities of conditional
densities for almost every endpoint.  For completeness, choose endpoint
transformation parameters in a small chart.  Fubini gives a base endpoint
for which the identities hold for almost every parameter.  The
transformed integral is smooth by dominated differentiation, using
\eqref{eq:init-conditioned-tail} and the Cameron--Martin exponential.
It consequently supplies a smooth version throughout that chart.
The versions agree on overlaps.  This proves the derivative statement
for the unnormalized conditional measures, including on closed profile
derivative supports.
\end{proof}

Choose a fixed large constant $D_s$ and let $\zeta_s$ be a smooth profile
which is one for $\mathcal N\le D_sp$ and zero for
$\mathcal N>2D_sp$.  We next compute the conditional integral on this
profile.  Write
\[
 X(s)=Z(s)+sa,
\]
where $Z$ is a standard Gaussian bridge and $a$ an independent standard
four-dimensional Gaussian.  Let $D$ be the vector of prescribed small
endpoint displacements: its quaternion part is the principal left
logarithm and its fourth component is the actual, unscaled $v$
increment.  Put $A=\sqrt r\,g a$.

\begin{lemma}[Endpoint disintegration on the profile]
\label{lem:init-endpoint}
On the profile $\zeta_s$, and for endpoints through the chart extension,
the endpoint equation has the unique small solution
\[
 A=A(D,gZ)=D+\text{terms of total degree at least two}.
\]
Its coefficients and its fixed normalized derivatives satisfy the
analytic path bounds of Lemma~\ref{lem:init-path-estimates}.  With respect
to bridge Gaussian measure, the conditional row density is the profile
with this solution substituted, multiplied by
\begin{equation}
 C_g e^{-|D|^2/(2rg^2)}
 \exp\!\left\{-\frac{|A|^2-|D|^2}{2rg^2}\right\}
 \frac{J(0)}{J(D_q)}\det(\partial_D A).
 \label{eq:R12}
\end{equation}
Here $J$ is the Haar Jacobian in quaternion logarithm coordinates, and
$C_g$ is $g^{-4}$ times a smooth positive function of $r$ and $R$.
\end{lemma}

\begin{proof}
The first-order tangent of the developed endpoint is $A$, because
$Z(1)=0$.  The iterated-integral bounds show that its derivative with
respect to $A$ is uniformly close to the identity when $gp$ is small.
The inverse-function estimate applies on a ball containing the entire
profile.  In that profile $|a|\le Cp$ and
$\mathcal N(Z)\le CD_sp$, so every possible solution lies in this
ball and the inverse is unique there.  One may use a fixed larger ball
of the same form for the prepared smooth extensions.

Changing the Gaussian endpoint variable $a$ to $D$ gives the determinant
$\det(\partial_D A)$ and the exponential $e^{-|A|^2/(2rg^2)}$.
Converting logarithmic coordinates to Haar endpoint density divides by
$J(D_q)$.  Separating its value at zero and the leading Gaussian gives
\eqref{eq:R12}.  The scale and coordinate-normalization factors are
exactly $C_g$.  The determinant is positive on the chosen inverse chart.
All derivatives follow from the analytic inverse and
Lemma~\ref{lem:init-path-estimates}.

Reversal can be imposed on this rule: reverse $X$ with its origin
subtracted, reverse the bridge, and use the averaged vertical convention
already specified.  The dyadic size test is invariant under reversal.
Thus neither the disintegration nor its extension breaks a required
symmetry.
\end{proof}

\subsection{Regular rows and the strip Gaussian}

We now partition the exact strip integral.  Its primary events are
endpoint flags and, on rows eligible for smooth calculation, selection
of the complement $1-\zeta_s$.  Delete a strip row whenever its horizontal
endpoint bond is flagged or it touches a vertical endpoint flag at
either column.  Such a row needs no profile split.  On all other rows
insert $1=\zeta_s+(1-\zeta_s)$, deleting the row in the second term.
The rows not deleted in a given term are called \emph{regular rows}.
Every pair of adjacent regular rows has the horizontal and vertical
bond charts, with consistent integers on their strip cell.  The
principal vertical endpoint constraint, including its tie convention,
is imposed once; it is automatic on these charts.

Primary events are grouped into the connected footprints used in
Section~\ref{sec:rg}.  Their supports include all positive and negative
flag and pattern queries.  Every interaction touching a deleted row
will belong to the corresponding compulsory factor.  Here
``compulsory'' describes the fixed-pattern integration: the factor must
retain the selected primary event and all its required reads.  It need
not have nonempty endpoint inventory, since a stochastic deletion alone
can occur with every endpoint bond unflagged.  After integration,
connected components with nonempty endpoint inventory remain in the
mandatory covering sum; components with empty endpoint inventory enter
the background-log reassembly.  These are separate decisions.

On regular rows use \eqref{eq:R12}.  Besides the horizontal term
$|D|^2/(2rg^2)$, the leading quadratic energy on a bond of the regular
vertical array is
\begin{equation}
 \frac{r}{2g^2}\int_0^1
 \left|J_s(D^0,D^1)+\sqrt r\,g\,dZ(s)\right|^2\,ds.
 \label{eq:R13}
\end{equation}
Here $J_s(c,c')=(1-s)c+sc'$, $D^0,D^1$ are the small vertical
logarithm/increment vectors at the two columns, and $d$ is the incidence
difference on the array of regular rows.  In particular, $dZ(s)$ is a
difference between neighboring bridges at time $s$, not a stochastic
differential.  The endpoint vectors are treated as bond data; they are
not asserted to be an exact linear field gradient.
In local relative-site charts they agree with that gradient at first
order.  The difference begins in degree two and consequently affects
the interaction exponent only at cubic order.  On the profiles, the
supremum and normalized derivative bounds for the difference between
the exact interaction and \eqref{eq:R13} are
$Cg\operatorname{poly}(M,p)$.  The two correction factors in
\eqref{eq:R12} have the same positive normalized order.  Fixed Taylor
remainders have their corresponding powers of $g$ times these fixed
logarithmic costs.

The length of the vertical array may grow with the rectangle.  It is
therefore necessary to perform the Gaussian completion with bounds per
row, not with a constant depending on the number of rows.

\begin{lemma}[Uniform strip completion]
\label{lem:init-completion}
For a finite periodic regular-row array or an array with arbitrary deleted
rows, the Gaussian completion of \eqref{eq:R13} admits finite-range
numerical means and covariances.  The covariances have uniform positive
upper and lower bounds.  The means and covariance deviations from the
identity have exponential position bounds and summable estimates over
bridge modes.  Their exact Gaussian ratio factors satisfy the joint
moment and support estimates used in Lemma~\ref{lem:rg-joint}.
The stationary energy is the strip energy \eqref{eq:R3}, up to local
errors $C\operatorname{poly}(M,p)e^{-c_LM}$ in chart norms.  These
assertions are uniform in array length and in the deletion pattern.
\end{lemma}

\begin{proof}
We first choose a Gaussian law with finite-range covariance and retain
its exact density ratio.  We then identify its stationary energy and
scalar determinant, and verify that profile failures still have a
Gaussian reserve under this law.

\emph{The numerical Gaussian law.}
In normalized sine coefficients write $Z_m=\xi_m/(\pi m)$; the
$\xi_m$ are independent white normals, with four colors at every regular
row.  The precision and linear term of the completed energy are
\[
 A_m=I+\frac{r^2}{(\pi m)^2}d^*d,
 \qquad |(L_m)_y|\le \frac{Cp}{m^2}.
\]
The bound on $L_m$ follows because the sine coefficient of the affine
interpolation is $O(m^{-1})$.  Shrink the fixed neighborhood of $r=1$
so that $4r^2/\pi^2<\delta<1$ there.  Approximate $A_m^{-1}$ by the
Neumann polynomial $C_{m,S}$ cut off at a sufficiently small fixed
multiple $S$ of $M$.  On each real spectral value $x\in[0,\delta/m^2]$,
\[
 C_{m,S}=\sum_{j=0}^{S}(-x)^j,
 \qquad c\le C_{m,S}\le1,
 \qquad |A_mC_{m,S}-I|\le (\delta/m^2)^{S+1}.
\]
These scalar estimates give the operator bounds also on arrays with
cuts.  In absolute row sums the same Neumann series is dominated by
$\delta/m^2$.  It gives exponential localization in the vertical
position and a summable $m^{-2}$ deviation from identity per row.

Set $\mu_m=-C_{m,S}L_m$ and $\eta_m=\xi_m-\mu_m$.  Use the Gaussian
probability with covariance $C_{m,S}$ and mean zero for $\eta_m$.
This change is exact when accompanied by the factor
$\det(C_{m,S})^{1/2}$ per color and the exponential with exponent
\[
 \frac12\eta_m^*(C_{m,S}^{-1}-A_m)\eta_m
       -(A_m\mu_m+L_m)^*\eta_m.
\]
The shift of the mean contributes
\[
 \frac12\mu_m^*A_m\mu_m+L_m^*\mu_m.
\]
Adding this contribution to the unchanged endpoint quadratic terms in
\eqref{eq:R13} and the horizontal energy gives the numerical stationary
energy.  The determinant, stationary energy and ratio exponent together
make the change of Gaussian law exact.

These changes may first be performed with finitely many bridge modes.
The covariance perturbations have summable trace per site and the means
are square summable, so the Gaussian laws and their density formulas
have limits as the number of modes tends to infinity.  The limiting
path laws are locally absolutely continuous with respect to independent
bridges.

\emph{Local factors from the exact ratio.}
We detail the locality of these corrections.  A residual
$A_mC_{m,S}-I$ has bandwidth $S+1$ and size
$(\delta/m^2)^{S+1}$.  Expand inverse and determinant ratios in this
residual, inserting a local $A_m$ where needed.  A chain of $h$ hops can
be grouped with its whole interval window and assigned the bound
\[
 m^{-2}(Ce^{-c_LM})^h,
\]
with fixed $M,p$ powers for its arguments.  There are only
$\operatorname{poly}(M)$ choices at each hop, and its full window has
load $C(h+1)$.  Group all modes of a given window in one factor.  The
resulting exponential envelopes involve
\[
 1+\sum_{m\ge1}m^{-2}|\eta_m(y)|^2
\]
at their read endpoints.  More explicitly, for a grouped window chain
$\gamma$ of $h$ hops we may take a single envelope coefficient
$e_\gamma\le\operatorname{poly}(M,p)(Ce^{-c_LM})^h$; the mode weight
$m^{-2}$ stays inside the displayed quadratic envelope.  Its
eighth-root support-hit sum is bounded by a sum of the form
\[
 \sum_{h\ge1}\operatorname{poly}(M,p,h)
       (C M^D)^h e^{AC(h+1)}e^{-c_LMh/8},
\]
which is small for large $H_N$.  In particular, no eighth root is taken
separately on the summable mode weights.  The determinant estimate for
products of exponential-minus-one factors now applies first in finitely
many modes.  Its small per-site endpoint load and trace hypotheses hold
uniformly because $\sum_m m^{-2}<\infty$.  Passage to all modes gives
the same joint bound, including the required eighth-root summability.
There is no factor proportional to the number of bridge modes.

\emph{The scalar and stationary energy.}
For the determinant of $A_m$ itself, group its convergent log series by
length and diagonal anchor.  Terms below the numerical cutoff are
extracted locally; beyond it they have the exponential window prices
just established.  Use the value on the array without holes for the
bulk scalar, and retain the local changes near holes in the compulsory
factors.  The exact full completion of \eqref{eq:R13} gives
\eqref{eq:R3}.  Numerical stationary coefficients and the prescribed
polynomially truncated square-root rows differ from it by exponentially
small weighted row sums, also on cut arrays.  Assign each coefficient
difference its endpoint and stencil reads.  This gives the claimed
local error in chart norms.

\emph{Profile failures under the numerical law.}
The numerical mean paths have
$\mathcal N$ size $Cp$ and bounded first derivative of that size, since
the mean sine terms and their first derivatives are absolutely summable.
For a specified set of rows, the centered numerical Gaussian marginal
has density bounded above by $\exp(C\,\#\text{rows})$ relative to the
independent bridge marginal: its covariance is at most identity, and
the lower and trace bounds control its determinant.  Consequently,
a failure of a regular-row profile forces the centered path to have
size at least $cD_sp$, also on derivative supports, once $D_s$ is large.
Here the mean bound, chart extension through $4t_*$, and the inverse
in Lemma~\ref{lem:init-endpoint} are all used.  Lemma~\ref{lem:init-path-estimates}
therefore gives the shared Gaussian-tail estimate required in the
joint bound.  Extend ordinary exponent vertices with larger smooth
path-size profiles before opening their failure factors.  Outside
that extension the good profile is zero, so the failure has the same
reserve.  Profiles attached to compulsory factors are retained on all
their regular-row reads.
\end{proof}

\subsection{Comparison with the retained energy near deleted rows}

The strip completion has now supplied the correct Gaussian and local
Taylor factors.  To obtain the retained class, we must still pay for
all extracted kinetic rows near deleted paths.  This is where the true
endpoint flags and the stochastic path deletions have different roles.

Away from their footprints, compare the completed squares with the
numerical rows of the full strip kernel without holes.  In a common
local quaternion chart, replacing the logarithmic bond vectors by the
transported chords changes a row by
$O(\operatorname{poly}(M)t_*^2)$.  The circle entries already agree
after the scale $R$ is included.  These are the positive-order Taylor
factors described above.  The exponentially accurate affine
normalization of the numerical rows is included among their local
errors.

Near primary events we use the following ownership rule.  Include
numerical and Taylor rows, and stationary and determinant compensations,
through $8M$ of a primary footprint; include defaults through $14M$,
with numerical ranges chosen inside the available margins.  Assign each
stationary squared row and its base extraction row to the same location.
For a vertical direct row, half lies in each adjacent strip.  Every
location is assigned exactly once, including locations belonging to a
compulsory factor.  Assign every deleted row measure, divided by the
same $C_g$ as in \eqref{eq:R12}, and every interaction touching that
row, to its compulsory component.  If an interaction touches several
deletions their joining rule gives a unique component.

A near-core vertex involving a regular row retains its quaternion chart
and actual path-size profile.  It may use the larger smooth extension
for its optional Taylor factor, but that does not remove its actual
profile.  Such vertices are exponentiated inside the compulsory factor.
Keep each actual regular-row profile once: through $14M$ it is retained
there, and farther away its failure can be opened as in
Lemma~\ref{lem:init-completion}.  Thus an interaction touching a deleted
path always retains the profile of any regular path it reads.  Every
factor depending on the regular-row pattern, or on the nearest surviving
bonds, retains those pattern queries in its support.  These are the original flag and deletion queries fixed in the initial
partition.  A later expansion of a retained regular-row profile as
$\zeta_s=1+(\zeta_s-1)$ does not change which rows the numerical
Gaussian regards as regular: $d$, $C_{m,S}$, and the stationary energy
continue to use that fixed primary pattern.  These rules specify the
integrand exactly and avoid counting an additive energy row twice.

\begin{lemma}[Initial compulsory-factor reserve]
\label{lem:init-reserve}
Choose $c_0,d_0$ sufficiently small, then the no-flag radius $D_0$
sufficiently large, and finally the path-profile constant $D_s$
sufficiently large.  These choices precede the choice of $L$.
For every fresh compulsory component, using the bounded-load record
specified below and denoting its load by $s_c$, its integrated factor
has the bound
\[
 \exp\{-cp^2s_c+C_LM^D(1+|\log g|)s_c\},
\]
or a stronger bound, with a positive reserve in the flagged increments.
The bound is conditional on all regular-row Gaussian arguments read by
the factor and holds on the differentiated profile supports as well.
It also holds on the no-flag extensions used for empty-output factors.
\end{lemma}

\begin{proof}
Let $n_c$ be the number of primary events in the component.  Choose its
fresh record with
\[
 s_c\le C n_c.
\]
Indeed, each primary carries only fixed-radius halos in units of $M$,
and a spanning tree of the proximity links joins their footprints by
paths of length $O(M)$ per link.  Every halo and joining path therefore
has bounded load.  We use this particular record, without arbitrarily
adding unused load.  Consequently a reserve $cp^2$ per primary gives
$cp^2s_c$ after decreasing $c$.

First consider only the deleted integrations and their touching
interactions, with the normalization $C_g$ removed.  A horizontal flag
has the reserve $cb r_e^2$ from \eqref{eq:init-heat-bound}.  A stochastic
deletion has the reserve $cD_s^2p^2$ from
\eqref{eq:init-conditioned-tail}.  For a vertical flag both its row
paths are deleted.  The real interaction loss is
\[
 1-\cos U(s)\,q(s)\cdot q'(s),
\]
one half of the squared normalized matrix-spin chord.  At the flagged
endpoint the principal vertical branch makes this at least $cr_e^2$.
The triangle inequality for the matrix-spin chord gives the lower bound
\[
 1-\cos U(s)\,q(s)\cdot q'(s)
 \ge cr_e^2-Cg^2(\mathcal N_1^2+\mathcal N_2^2),
\]
using the path displacement estimate and decreasing $c$.
Use only a sufficiently small fixed fraction of the interaction loss
for this payment.  Its adverse term then has an exponential moment
controlled by Lemma~\ref{lem:init-path-estimates}.  A fixed H\"older
inequality retains positive fractions of the horizontal and stochastic
reserves.  Each row touches only boundedly many bonds, and if both
endpoints are flagged their payments share the same loss.  All other
touching interactions have modulus at most one.  Endpoint derivatives
are treated by the last part of Lemma~\ref{lem:init-path-estimates};
their fixed powers of $g^{-1},M,p$ are charged to the displayed entropy
term.

The direct extracted energy \eqref{eq:R4} uses only a small fraction of
these reserves.  At flags it is linear in the large increment, while
the payment is quadratic.  At nonflag bonds incident on a deletion,
the direct squares absent from the regular array cost at most
$Cc_0bt_*^2$ per bond.  Their number is a fixed multiple of the number
of primary events, and the vertical direct energy has already been
split between its two strips.  Choosing $c_0,d_0$ small leaves a positive
reserve.  Chord/logarithm differences on included direct rows have the
Taylor bounds already established.

It remains to pay for the masked second-stack rows.  A row whose mask
is off extracts no energy.  For a row whose mask is on, both the base
row and the regular-array row have size $Ct_*$.  If its center has
distance $a$ from the nearest primary, their difference, in a common
local frame, is bounded by
\begin{equation}
 C\operatorname{poly}(M)t_*^2+Ct_*e^{-ca}+Ct_*e^{-c_LM}.
 \label{eq:init-row-difference}
\end{equation}
To prove this, use a flat chart through a fixed fraction of that
distance, within the stencil.  The chord/logarithm difference gives
the first term.  The kernels on arrays with and without holes agree
up to their exponential tails at that distance, giving the second;
numerical truncation gives the third.  Farther increments of a base
row can be large, but its mask permits growth only with the prescribed
exponential profile, slower than the row decay when $R_*$ is chosen
large.  Regular-array increments are all chart sized.  For small $a$
the separate row bounds give the same estimate.  Beyond the assigned
core ranges this distance comparison is unnecessary because the
ordinary local factors are used instead.

Using $\bigl||a|^2-|b|^2\bigr|\le(|a|+|b|)|a-b|$ and summing the
distance term in \eqref{eq:init-row-difference} costs at most
$Cbt_*^2e^{-c'D_0}$ per flag: its masked-on rows begin at distance
$D_0$.  It costs at most $Cbt_*^2$ per stochastic primary.  Choose $D_0$
large and then $D_s$ large to pay these costs with their respective
reserves.  On a no-flag extension only stochastic primaries occur, so
the same argument applies.  Finally, the Taylor and numerical errors
near a component cost at most
\[
 CbM^D(t_*^3+t_*^2e^{-c_LM})
\]
per primary.  Since $t_*\asymp gp$, these leave a positive fraction of
$p^2$ after the stated parameter choices.  The support and integration
costs are bounded by the displayed $C_LM^D(1+|\log g|)s_c$.
\end{proof}

The margins in this construction have the same meaning as in the exact
step.  Each flag creates only boundedly many deletion sites; the original
interactions connect neighboring rows; and the numerical covariances
have finite range.  Long inverse or determinant chains retain every
successive window in their complete supports.  Therefore disjoint
complete supports have zero cross covariance and disjoint deleted-path
variables.  Their marginals do not change when remote primary events
are removed, and their integrals factor exactly.  Overlapping supports
are controlled by the joint estimate, not by an independence assertion.
This is precisely the distinction used in
\cite[Section ``Connected sums and exact reassembly'']{SharpO4}.

\subsection{Canonical coefficients and the initial scalar}

We can now apply the joint product, rooted-tree and background-log
estimates of Section~\ref{sec:rg}.  The required inputs have just been
verified: ordinary vertices have positive-order deterministic bounds
on their extensions; ratio factors have the tilted determinant bounds
of Lemma~\ref{lem:init-completion}; failures and compulsory factors have
the shared reserves of Lemma~\ref{lem:init-reserve}; and complete supports
give exact factorization for disjoint collections.  Fixed derivative
orders incur only the polynomial costs already included in the choice
of $P_0$.  The background logarithm is taken only for empty-output
polymers, with the inventory-preserving corrections of the exact step.
Thus this operation reconstructs the original path integral exactly.

For clarity we verify that the continuous paths also give the canonical
path norms required by Lemma~\ref{lem:rg-canonical}.  Strengthen the masks
for terms through normalized order four and use full kernels in their
Taylor coefficients.  Each Taylor coefficient of \eqref{eq:R12} and of
the interaction is a finite Gaussian-chaos polynomial in iterated
integrals, including deterministic interpolation and mean paths.
Dyadic polygonal approximations converge in every fixed moment, by the
rough-size estimates of Lemma~\ref{lem:init-path-estimates}, which hold
uniformly for those approximations.

In the full completion write the centered bridge modes as
$A_m^{-1/2}$ times independent white row modes.  The deviation of this
matrix from the identity has size $Cm^{-2}$ with exponential weight in
the row separation.  Its contribution to the paths is consequently a
$C^1$ Gaussian path with coefficient moments exponentially summable in
row origin; the differentiated sine series is absolutely summable at
that order.  Mean substitutions have the same property.  Expanding a
fixed iterated polynomial after these substitutions preserves exponential
position summability by the smooth-drift rough-path bounds.  A connected
expectation of $v$ vertices must share independent-row Gaussian variables
between them, and a surviving Wick term has at least $v-1$ intervertex
pairs.  This is immediate on polygonal approximations and passes to
each homogeneous Taylor coefficient.  It is exactly the Gaussian
count used in the canonical calculation of Section~\ref{sec:rg}.

The same square-root power-series estimates compare numerical and full
covariances with exponential accuracy in $M$, retaining their row
positions.  The preceding moment argument therefore proves the
numerical-to-full replacement with the canonical path weights.  Lift
all paths before folding them into a torus.  A path that winds retains
its length and is assigned the same winding price as in
Section~\ref{sec:rg}; no winding coefficient is discarded.  Taylor
remainders on the guards, and removal of the guards using the fixed
polynomial moments, have the required order-five estimates.  There is
no old regular-error transfer in this initial integration.

Extract the constant, linear and directional quadratic terms through
order four.  Absorb the imaginary linear terms into the phase.  Use
the corrected couplings as the starting shapes of \eqref{eq:R3}, making
exactly the shape-change and displayed-power Taylor adjustments specified
with \eqref{eq:R6}.  The harmonic strip parameters need agree with the
starting shape only before this correction.  The remaining canonical
possibilities are exactly \eqref{eq:R5}.  Apply the regular-error,
covering and winding normalizations already proved for the exact step.
The error in each displayed affine correction is
$O_L(H_N^{-1.05})$, and the finite coefficient bounds fix the initial
canonical caps.  This proves all class assertions of
Proposition~\ref{prop:init-density}.

The scalar has a particularly simple leading form.  The extracted
interaction contributes $br$ per site and the endpoint heat density
contributes $-4\log g$.  The order-zero Gaussian determinant is a bounded
Lipschitz function of $r,R$.  The first-order scalar term vanishes by
Gaussian parity; scalar corrections start at order two.  Hence
\begin{equation}
 V_B^{\mathrm{init}}
   =br-4\log g+V_{\mathrm{init}}(r,R)+O_L(1/b),
 \label{eq:R14}
\end{equation}
where $V_{\mathrm{init}}$ is bounded and Lipschitz on the parameter band.
The normalization $C_g$ is the same on regular and deleted rows, so this
is one bulk scalar, independent of the periods and of the time seam.

\section{Matching the regulators and integrating the comparison}
\label{sec:comparison}

We now have two exact retained-density evolutions.  Their local
coefficients contract when their four affine couplings agree, but those
couplings must first be matched.  We choose the four initial parameters
$(b,r,c_x,c_t)$ by a topological shooting argument.  We then integrate
the resulting small differences.  The latter step uses the positivity
of the flat reference, because a covering sum need not admit a useful
pointwise relative bound.
The resulting partition comparison proves
Proposition~\ref{prop:rg-comparison}.  We also use the extracted scalars
to bound the pressure under a small fugacity change; this will locate
the particle sectors needed in the final mass calculation.

\subsection{Intrinsic affine coefficients and continuity}

The Wilson reference evolution must first be identified with the
trajectory whose depth was calibrated in \cite{SharpO4}.

\begin{lemma}[Reference evolution and parameter continuity]
\label{lem:cmp-continuity}
On the flat reference, the quaternion affine couplings are exactly those
of the nearest-neighbor evolution in \cite{SharpO4}, and the physical
circle affine stiffnesses, expressed in the original $v$ coordinate,
do not change under blocking.  At fixed $H,N$ and fixed scale factors
$R_j$, all four displayed couplings of either evolution are continuous
functions of the initial parameters along every strictly admitted finite
segment.
\end{lemma}

\begin{proof}
On flat cochains the observation preserves the lifts and decouples the
quaternion and real-field integrations.  To see that the displayed
affine coefficients are independent of the particular retained-density
representation, fix an untwisted output near the identity with integers
zero.  Perturb a single quaternion Cartan direction, or the real field,
by sine and cosine waves in one lattice direction.  On sufficiently
small perturbations every mask is one and the covering sum is one.
The phase sum is constant.  Divide the second variation of the negative log
density by volume and by the mean square bare wave gradient, averaging
the sine and cosine versions.

As the wave momentum $p_{\mathrm{mom}}$ tends to zero, this quotient
converges to the corresponding displayed affine coefficient.  Indeed,
at an anchor subtract the common field value.  The remaining wave is
its affine part plus a second-difference remainder.  In the weighted
direction norm, these have sizes respectively
$C_L|p_{\mathrm{mom}}|/t$ and $C_L|p_{\mathrm{mom}}|^2/t$.
The canonical compensations and regular errors have zero affine
quadratic jet, so their contribution to the quotient tends to zero.
Their anchored derivative sums converge by their defining norms.
Winding errors at the lowest nonzero momenta have exponentially small
period weights, and their direction norms require only the wave bound
itself.  The kinetic rows have their exact affine normalization and
uniform row moments.  In the quaternion calculation there is no
acceleration term from subtracting the anchor rotation, because the
wave remains in one Cartan direction.  Increasing dyadic square tori
suffice, and the convergence is uniform in the retained class at each
fixed layer.

For the flat reference at locally trivial output winding, fixing the
output variables and integers fixes the common integration sheet.
The reverse-tree parametrization of the observation leaves a genuine
real Gaussian integration, with no additional periodization of its
zero mode.  The local output log density is therefore the quaternion
endpoint log of \cite{SharpO4} plus the log of the observed real Gaussian
field.  The quaternion observation, including its fallback branch, is
exactly that of the cited evolution; induction identifies its
coefficients.  Gaussian observation preserves affine stiffness in the
physical coordinate.  The intrinsic characterization just proved shows
that different cutoff representations do not affect either conclusion.

For continuity, first work on a fixed finite torus and at fixed output
arguments.  Brownian initial densities vary continuously with the heat
time: use heat-bridge continuity, or finite time-slicing and the heat
bounds of Lemma~\ref{lem:init-path-estimates}.  The potentials are
continuous and deck sums have Gaussian domination.  Sheet predicates
are held fixed.  The block kernels preserve continuity by dominated
convergence; their geometric and medoid predicates do not depend on
the varying coupling parameters.  The direct-energy and covering bounds
provide domination.  On the all-good domain the retained representation
is a nonzero exponential.  Its finite-period log ratios have uniformly
bounded higher derivatives, since the possibly large global phase is
constant there.  Finite-difference approximations to their second variations have a
uniform error by these higher-derivative bounds; hence those second
variations are continuous.  Taking the uniform small-momentum limit identifies the next affine
coefficient continuously.  Induction proves the statement at every
finite admitted depth.  No bound on a derivative of the total accumulated
phase is needed.
\end{proof}

For the prescribed Wilson coupling $\beta$, use the depth supplied by
\cite[Proposition ``Admission and shooting'']{SharpO4}.  By our definition
$R_j^2=\beta/H_j$, and the cited proof gives a quaternion running coupling
with ratio to $H_j$ close to one.  With $H$ large, the flat reference thus lies strictly
inside the shape, stiffness-ratio and coupling bands of
Section~\ref{sec:rg}, uniformly along the entire trajectory.

\subsection{Four-parameter shooting}

\begin{proposition}[Matched trajectories]
\label{prop:cmp-shooting}
For every sufficiently large $\beta$ and the depth just specified, there
are parameters
\[
 b=\beta+O_L(1),\qquad r=1+O_L(\beta^{-1}),\qquad
 c_x,c_t=\beta+O_L(1)
\]
for which both evolutions are admitted through layer zero and, with the
distances of \eqref{eq:R7},
\begin{equation}
 |\lambda_j|+U_j\le C_L\rho^{N-j},\qquad \lambda_0=0,
 \qquad \max(q_*,L^{-1})<\rho<1.
 \label{eq:R15}
\end{equation}
The constants are independent of $\beta$ and of the admitted periods.
\end{proposition}

\begin{proof}
Parameterize the initial values by the four uncorrected discrepancies
$x=(x^q,x^w)$:
\[
 x^q=b(r^{-1},r)-(\beta,\beta),\qquad
 x^w=R_N^{-2}\{b(r^{-1},r)-(c_x,c_t)\}.
\]
Every fixed ball in these coordinates is available for sufficiently
large $\beta$.  Indeed the first two coordinates determine $b$ and $r$
by product and ratio, and the last two then determine $c_x,c_t$.
Proposition~\ref{prop:init-density} gives
$|\lambda_N-x|\le C_L$, while all initial distance terms are uniformly
bounded.

Choose $\rho$ as stated and let
\[
 z_h=\rho^{-h}\lambda_{N-h},\qquad 0\le h\le N.
\]
Run each shot until its first possible exit from $|z_h|\le D$.
Convolution of the first estimate in \eqref{eq:R7}, followed by the
second, gives on every such segment, including its first attempted
output,
\begin{align*}
 U_{N-h}&\le (C'_L+C'_LH^{-c}D)\rho^h,\\
 |z_{h+1}-\rho^{-1}z_h|&\le C''_L+C''_LH^{-c}D.
\end{align*}
Choose $D$ large and then $H$ large so that the right side of the second
inequality is much smaller than $(1-\rho)D$.  The comparison with the
reference then puts the entire barrier, and the next coefficient step,
strictly inside the admitted bands.

We give the topological argument explicitly.  Use the initial $x$-ball
of radius $D+C_L+1$.  Suppose no admitted shot in this ball has zero
terminal discrepancy.  If its initial discrepancy is outside the
barrier, map the shot to its radial direction.  Otherwise join successive
$z_h$ by line segments and use the direction at its first crossing of
$|z|=D$.  If it has no crossing, use the direction of its nonzero
terminal value.  This defines a map from the whole parameter ball to
the unit sphere.

The map is continuous.  On $|z_h|=D$ the next segment points strictly
outward already at its start: the scalar product of $z_h$ with
$z_{h+1}-z_h$ is positive by the preceding estimate.  Thus crossings
persist under perturbation and no discontinuity occurs when a first
crossing moves across a discrete layer.  At a terminal crossing the
crossing and terminal-direction definitions agree.  Parameter
continuity is Lemma~\ref{lem:cmp-continuity}.  On the boundary of the
initial ball the estimate $|\lambda_N-x|\le C_L$ makes this map
homotopic to the radial identity.  A map of nonzero boundary degree
cannot extend over the ball, a contradiction.  A zero terminal
shot therefore exists and never exits its barrier.  The two estimates
above prove \eqref{eq:R15}.
\end{proof}

There is an important consequence of the exact equality $\lambda_0=0$.
After removing the bulk scalars and the propagated phases, the
\emph{unbounded} direct kinetic terms in the two terminal exponents
agree exactly.  Every remaining second-stack row is masked and has a
uniform size bound.  At positive free-history depth use identical
symmetric numerical truncation and affine-correction rules in both
runs.  The weighted free-kernel difference then retains its exponential
history gain.  Together with \eqref{eq:R15} this gives, on a terminal
rectangle of volume $V=m_xm_t$,
\begin{equation}
 \|X_B-X_F\|_\infty\le C_HV e^{-c_H\beta},
 \qquad
 \sup_y\sum_{\ell:y\in P_\ell}
 e^{2s_\ell}\|k_{\ell,B}-k_{\ell,F}\|_\infty
 \le C_He^{-c_H\beta}.
 \label{eq:cmp-terminal-distances}
\end{equation}
Here $X$ is the scalar- and phase-stripped log exponent, and the
covering weights use a common label list.  Polynomial factors from
the weighted distance are included in the exponential by decreasing
$c_H$.  The reference assignments can be made real with phase zero:
all integrations and logarithmic reassemblies then preserve reality
and ordinary $w$ negation.  In the paired construction, changing a
reference phase on a fine contribution already forced to vanish does
not change those assignments.  We are now ready to integrate
\eqref{eq:cmp-terminal-distances}.

\subsection{From matched coefficients to matched integrals}
\label{subsec:comparison-integration}

The terminal estimate \eqref{eq:cmp-terminal-distances} compares the exponents and individual
covering weights.  To deduce \eqref{eq:R1}, we must integrate that comparison
without dividing by a possibly small signed covering sum.  We use the
integrated argument of \cite[Section~7, ``Integrated comparison of endpoint
densities'']{O3Continuum}.  Its three inputs are an exponential moment of
selected kinetic rows, a measurable replacement of rough regions with an
energy gain, and the exact identities for mandatory covering sums.  We
verify these inputs for the flat lift, including its integer data, both
time seams, and rectangles of unbounded aspect ratio.

Fix the large terminal parameter \(H\), and put
\[
 p_0=(\log H)^{P_0},\qquad t_0=H^{-1/2}p_0,\qquad
 P=\pi R_0,\qquad V=m_xm_t.
\]
Here \(P\) is the length of the terminal circle coordinate; it is unrelated
to the fixed exponent \(P_0\).  The terminal torus is
\(\Lambda=(\mathbb Z/m_x\mathbb Z)\times
(\mathbb Z/m_t\mathbb Z)\).  Its periods are dyadic and exceed a constant
depending on \(H\).  We will increase this constant to accommodate the
meshes and tiles below.  No bound on \(m_x/m_t\) is imposed.

Write \(z=(q,w,k)\) for the retained variables, with
\(q_x\in S^3\), \(0\le w_x<P\), and one integer \(k_e\) for a chosen
orientation of each unoriented bond.  Integration means product normalized
Haar measure in \(q\),
Lebesgue measure in \(w\), and counting measure in \(k\).  Denote this
measure by \(d\nu\), and let \(\mathcal F=\{dk=0\}\) be its flat part.
The two time-seam prescriptions are denoted by \(\eta=+,-\).  They enter
the transported \(q\)-differences, while
\(u_e=dw(e)+Pk_e\).  After removing the volume scalar and the propagated
phase, the reference density on \(\mathcal F\) is
\(e^{X_F^\eta}\Xi_F^\eta\).  It is nonnegative for either seam.  Set
\[
 Z_{F,+}^{\rm end}=\int_{\mathcal F}e^{X_F^+}\Xi_F^+\,d\nu.
\]

For clarity, denote covering weights in this subsection by
\(\kappa_\ell\), reserving \(k_e\) for the edge integers.  Each label has
a complete site support \(P_\ell\), a load \(s_\ell\ge1\), and a fixed
nonempty bond inventory \(J_\ell\).  Its support contains the endpoints
of every integer argument as well as all positive and negative eligibility
tests.  The weights vanish unless \(J_\ell\) lies in
\[
 D(z)=\{e:r_e>t_0\},\qquad
 r_e^2=|\nabla_e q|^2+u_e^2.
\]
As in the retained class, compatible labels have disjoint complete
supports, and their inventories must partition \(D(z)\).

For a second mandatory covering sum \(\widetilde\Xi^\eta\), use a common
label list and set absent weights equal to zero.  Write
\[
 d_\ell=\|\widetilde\kappa_\ell-\kappa_{\ell,F}\|_\infty,
 \qquad
 b_\ell=\|\widetilde\kappa_\ell\|_\infty+
              \|\kappa_{\ell,F}\|_\infty.
\]
The sum \(\widetilde\Xi^\eta\) may have complex weights.  In the comparison
below, its combined weights satisfy the retained support-hit cap
\[
 \sup_x\sum_{\ell:x\in P_\ell}e^{64s_\ell}b_\ell
 \le C\exp(-p_0^{1/4}),
\]
where \(C\) is a fixed multiplier.  Our objective is an integral bound in
terms of \(\sum_\ell d_\ell e^{2s_\ell}\), uniform in the rectangle and
in \(P\).  We establish the analytic inputs before stating and proving
that comparison.


First we dispose of the part of the Brownian density outside \(\mathcal F\).  Every such
configuration requires a covering label retaining the curvature bound
\eqref{eq:R10}.  Summing that recurrence over the ultraviolet layers gives
\[
 F_N\ge c_L\beta^{3/2}(\log\beta)^{P_0}.
\]
After taking absolute values, discarding inventory and compatibility
constraints bounds all other covering sums by \(e^{C_HV}\).  The bounded
terms of the exponent have the same cost.  The remaining direct energy
controls the real coordinate and the integer sums: integrate along a
spanning tree, and use a reserved exponential in the increment to sum
each remaining edge integer.  This gives
\begin{equation}
 \int_{\mathcal F^c}|\text{Brownian endpoint density}|\,d\nu
 \le P\exp\{-F_N+C_HV\}.
 \label{eq:comparison-nonflat}
\end{equation}
The factor \(P\) is the unconstrained common translation.  Conversely,
an aligned cap in the trivial winding sector gives
\begin{equation}
 Z_{F,+}^{\rm end}\ge P\exp(-C_LV\log H).
 \label{eq:comparison-cap-lower}
\end{equation}
The cap radius is a fixed multiple of \(H^{-1/2}\); there are no flags,
the covering sum equals one, and its product measure and kinetic cost
give the displayed bound.  For the rectangles in \eqref{eq:R1},
\(V\asymp\beta^{3/2}(\log\beta)^6\).  Since \(P_0>20\),
\eqref{eq:comparison-nonflat} is negligible relative to
\eqref{eq:comparison-cap-lower}, even after an arbitrary phase is inserted.

\subsection{Kinetic-row moments for the flat reference}
\label{subsec:comparison-row-moments}

Group the nonnegative direct and second-stack kinetic terms at their
unoriented bonds.  Truncation at a sufficiently large multiple
\(D_1\asymp C_L(\log H)^2\), with all mask reads included in the stencil,
gives nonnegative rows \(Y_e\) and a decomposition
\begin{equation}
 X_F^\eta=-\sum_eY_e+X_{\rm rem},\qquad
 |X_{\rm rem}|\le C_LV,
 \qquad
 |X_{\rm rem}(z)-X_{\rm rem}(z')|\le C_Ln
 \label{eq:comparison-remainder}
\end{equation}
when the configurations differ at \(n\) sites and the incident integer
data.  Exponential spatial localization of the kinetic kernels, in both
row and column sums, controls the truncation tails.  The
regular and canonical fields give the support-hit bound in
\eqref{eq:comparison-remainder}; the additional cubic slot has per-anchor
size at most \(C_LHt_0^3=C_LH^{-1/2}p_0^3\), which still tends to zero
after multiplication by every fixed logarithmic factor needed here.

\begin{lemma}[Selected kinetic rows]
\label{lem:comparison-row-moments}
For a sufficiently large dyadic integer \(B_T\) comparable to a fixed
power of \(\log H\), suppose the two periods are multiples of \(2B_T\).  There exists \(0<\vartheta<1\)
such that, for every set \(I\) of unoriented rows and either time seam,
\begin{equation}
 \frac{1}{Z_{F,+}^{\rm end}}
 \int_{\mathcal F}e^{X_F^\eta}\Xi_F^\eta
       \exp\left(\vartheta\sum_{e\in I}Y_e\right)d\nu
 \le \exp(C_L|I|B_T^2\log H).
 \label{eq:R16}
\end{equation}
The choices may make \(1-\vartheta\) smaller than any prescribed fixed
inverse power of \(\log H\).
\end{lemma}

\begin{proof}
We first check the reflection input to the lemma ``Exponential moments
of selected rows'' in \cite[Lemma~7.3 and Equation~(7.5), Section
``Integrated comparison of endpoint densities'']{O3Continuum}.
At the fine scale, cut the flat reference into
two half cylinders and choose lifts in each half.  Their transverse
integer windings must agree.  Along either separating seam, fix an
origin and express each boundary site's sheet by cumulative tangential
integers.  Flatness then says that every cross-link integer is the
difference of these boundary sheet displacements plus one common
integer.  The common integers on the two seams can be summed independently.

In the lifted variables, the kernel for one seam is consequently a
product Gaussian kernel and Wilson \(q\)-kernel, summed over a common
deck shift.  Each unsummed kernel is positive semidefinite:
the Gaussian is positive definite, and
\(e^{\beta q\cdot Q}\) is positive definite by its power series.
The common-shift sum preserves this property.  Indeed, let \(g\)
denote simultaneous translation by one sheet, including its sign on
\(q\), and let \(K\) be the unsummed kernel.  For a finite set of
boundary configurations \(z_i\), the matrices
\[
 \frac1{2M+1}\sum_{a,b=-M}^{M}
       K(g^az_i,g^bz_j)
\]
are Gram matrices.  Simultaneous-shift invariance and Gaussian decay
show that their limit is \(\sum_{n\in\mathbb Z}K(z_i,g^nz_j)\).
The restrictions on the flat halves and their equal transverse winding
are diagonal restrictions on these positive forms.  Tangential
directions are chosen the same in the two halves, and normal integers
transform with their orientation.  This proves seam reflection
positivity of the untwisted flat reference.

A minus-center twist inserts the isometry \(q\mapsto-q\) in the seam
kernel.  Cauchy--Schwarz therefore bounds a twisted pairing by the
geometric mean of its two untwisted doubled-half pairings.  The twist
can be moved to any desired seam by a change of representatives.
These statements pass through the exact observations.  A half-observable
uses only internal output sites and integers; its fine data and
observation noises lie in that same half.  The observation kernels are
reflection equivariant on the lifts.  Cross-link output integers are
not read by such an observable and sum out by normalization.  Thus the
fine reflections can be taken at the block seams corresponding to the
terminal reflections.

We can now use the chessboard argument.  In a translated grid of
\(B_T\)-squares, retain the rows in \(I\) whose complete stencils are
internal to one tile, with a margin \(C_LD_1\).  The tile test is the
exponential of their sum.  Reflection disseminates a tile test to a
subset of the kinetic rows without repetitions.  To keep the integral
over the unbounded lift finite, use coefficient \(1-\delta\) in this
test, where \(\delta>0\) is a prescribed inverse power of \(\log H\).
The unspent fraction \(\delta\) of every direct term supplies spanning
tree integration and all deck sums.  With
\eqref{eq:comparison-remainder}, the absolute covering cap, and
\eqref{eq:comparison-cap-lower}, a fully disseminated test has normalized
integral at most \(e^{C_LV\log H}\).  Both its numerator bound and its
denominator lower bound have the same translation factor \(P\).

Taking the chessboard root gives a cost
\(e^{C_LB_T^2\log H}\) per occupied tile, and there are at most
\(|I|\) occupied tiles.  For a time twist, first place it at a tile seam
and apply the preceding Cauchy--Schwarz bound.  Both doubled arrays are
untwisted, so their geometric mean has the same tile-counting cost and
the same denominator \(Z_{F,+}^{\rm end}\).

Finally average over the \(B_T^2\) translations of the grid.  Each row
is internal for a fraction at least \(1-C_LD_1/B_T\) of them.
Generalized H\"older and \(Y_e\ge0\) give \eqref{eq:R16} with
\[
 \vartheta=(1-\delta)(1-C_LD_1/B_T).
\]
Reflection covariance permits the twist to be moved separately for
each grid.  Choose \(\delta\) small and then \(B_T\) large to obtain
the stated accuracy.  Truncating the nonnegative tests first and then
using monotone convergence justifies the argument for unbounded rows.
\end{proof}

\subsection{Replacing rough regions on a rectangle}
\label{subsec:comparison-edits}

The next construction is the flat-lift version of the lemma
``Replacement and its energy bounds'' in
\cite[Lemma~7.4 and Equations~(7.7)--(7.8), subsection
``Regions that can be replaced at an energy gain'']{O3Continuum}.
The all-scale marking below is the feature that removes the source's
bounded-aspect-ratio restriction.

For a flat configuration, lift the data to \(\mathbb Z^2\), with their
integer windings and prescribed sign twists, and put
\[
 f(x)=\max_{e\ni x}r_e,\qquad \alpha=\varepsilon t_0.
\]
The nonnegative function \(f\) is periodic, even though the lifted real
coordinate need not be.  Mark every site of every integer square on the
cover whose average of \(f\) exceeds \(\alpha\), including squares of
radius zero and of arbitrarily large radius.  Let \(U\) be the projection
of their union to \(\Lambda\).  Dilate \(U\) in max distance by
\(d=CD_1\), take the max-neighbor connected components, and retain the
complete components containing a bond of \(D(z)\).  Call these \(E_C\).
Their inventories \(D_C\) are the original flagged bonds in the corresponding
components; they form a partition of \(D(z)\).

\begin{lemma}[Flat replacements and their costs]
\label{lem:comparison-flat-edits}
The constants \(\varepsilon>0\) and \(C\) can be fixed, in that order,
so that the following holds for every sufficiently large \(H\) and every
admitted rectangle with both periods sufficiently large depending on
\(H\).  For every flat configuration and every selected family of its
regions \(E_C\), there is a measurable probability distribution of flat
replacement configurations that leaves all other sites unchanged, changes
only integers incident on replaced sites, removes exactly the inventories
of the selected regions, and creates no flagged bond.

If \(E\) is the union of the selected regions and \(n=|E|\), there is a
set \(I_E\) of rows, containing every changed kinetic stencil, such that
\( |I_E|\le C d^2n\).  For the original configuration \(z\) and every
sampled replacement \(z'\),
\begin{equation}
 \sum_{e\in I_E}Y_e(z)\ge cp_0^2n/d^2,
 \qquad
 \sum_{e\in I_E}Y_e(z')\le C_Ld^2p_0^2n.
 \label{eq:comparison-edit-energy}
\end{equation}
Relative to the site and incident-integer measure, the conditional fill
density is at most \(e^{C_Ln\log H}\) when \(E\ne\Lambda\), and at most
\(P^{-1}e^{C_Ln\log H}\) when \(E=\Lambda\).  Nonempty exterior data
retain their original winding.  In a whole-torus replacement one may
choose the winding, while keeping the prescribed time seam.
\end{lemma}

\begin{proof}
At every lift of a point of \(U^c\), all squares containing that point
have average gradient at most \(\alpha\).  The concentric-average proof
in the cited lemma applies on the whole cover: the discrete
\(L^1\) Poincar\'e inequality compares the point with averages on doubled
squares, at scale \(R\) at cost \(CR\) times the average gradient.
Compare two such averages at a scale comparable to their separation
and sum the geometric scales.  This gives, in the representatives
transported by \(k\),
\[
 |q_x-q_y|+|w_x-w_y|\le C\alpha|x-y|_\infty
 \qquad(x,y\text{ above }U^c).
\]
All sizes were admitted, so this estimate does not stop at the shorter
period of the rectangle.

Extend the real coordinate by distance cones.  For example, the infimum
of \(w_y+C\alpha|x-y|_\infty\) over prescribed lifts \(y\) is finite,
agrees with the prescribed data, and is Lipschitz.  The set of prescribed
lifts includes every deck translate, so the extension has exactly the
same additive winding under translation by a period.

For \(q\), first extend its four Euclidean coordinates near the
prescribed data.  Averaging the upper and lower distance-cone extensions,
and clipping to the bounded coordinate range if needed, makes this
construction equivariant under the required sign changes.  Within
distance \(c/\alpha\) of the prescribed set its norm is bounded away
from zero; normalization gives an \(S^3\)-valued extension with Lipschitz
constant \(C\alpha\).  Choose a rectangular mesh of spacing comparable
to \(c'/\alpha\), with \(c'\) sufficiently smaller than \(c\), and
preserve that extension in each mesh box meeting prescribed data.
Choose the remaining vertices and shortest edge arcs per mesh orbit,
then extend by the prescribed signs to all translates.  Each unfilled
box has a boundary map with a fixed Lipschitz bound after rescaling.

Such a loop has a uniformly Lipschitz filling in \(S^3\).  One concrete
construction chooses a point separated from the antipodal image of the
loop and normalizes the cone from that point to the loop.  The antipodal curve
has bounded length, so a fixed-radius tubular covering occupies less
than the volume of \(S^3\); a point outside it exists.  Ordered finite
nets and fixed measurable choices of shortest arcs give measurable
fillings.  This is the quantitative filling construction of
\cite[Lemma ``Quantitative filling on \(S^3\)'', Section ``Relative
comparison of the retained densities'']{SharpO4}.  Adjacent fillings use
the same previously fixed edge maps.  Defining choices by mesh orbits
also makes the resulting map sign-equivariant on the cover.  Both
periods exceed the mesh size, which is the only period-size requirement
in this construction.

Sample the replaced spins independently in balls of radius \(c't_0\)
around these extensions, using Haar measure in the sphere factor and
Lebesgue measure in the real factor.  Extend the samples to all
translates with the original winding, and reduce to the chosen
representatives, inducing the edge integers and the corresponding
signs on \(q\).  The extension equals the original data outside \(U\).
Taking \(\varepsilon\) and \(c'\) sufficiently small therefore makes
every new internal or boundary chord smaller than \(t_0\).  The collar
of each \(E_C\) ensures that no originally flagged bond crosses its
boundary.  Any subset of the regions can consequently be replaced
independently, with exactly its inventories removed.  Exterior and
unselected data keep their original representatives and integers.

If there are no prescribed data, perform the same mesh construction in
the trivial integer winding, with the fixed time sign still imposed.
For a whole-torus replacement one may use this construction and add a
uniform common translation to the real coordinate.  Each local sampling
ball has four-dimensional volume bounded below by a fixed multiple of
\(t_0^4\).  Hence the fill density is at most \(e^{C_Ln\log H}\).
In the whole-torus case the uniform translation supplies the additional
factor \(1/P\), with any remaining fixed or polynomial factor in \(H\)
included in the same exponential bound.

It remains to check the energy estimates without an aspect-ratio loss.
The image of a covering square may wrap around either period.  In a
wrapped coordinate the numbers of preimages of its sites differ by at
most a factor two; thus the multiplicities in a square's torus image
differ by at most a fixed factor.  A marked square therefore still
provides an average-gradient lower bound, with a decreased constant,
on its image.  Fixed dilations of these images have uniformly bounded
volume ratios.  Choose finitely many generating squares, possible
because the torus has finitely many sites, and select disjoint images
in decreasing order of size.  Fixed enlargements of the selected images
cover the generators.  Dilation by \(d\) costs at most \(Cd^2\), so their
total area in a region is at least \(c|E_C|/d^2\).

On each selected image, either increments greater than \(t_0\) provide
a fixed fraction of the average gradient, in which case the linear
branch of the direct energy gives the required lower bound, or the
quadratic branch does so by Cauchy--Schwarz.  In both cases the energy
is at least \(cH\alpha^2\) times its area.  Bonds in the selected
images have bounded overlap, and \(H\alpha^2=\varepsilon^2p_0^2\).
This proves the first bound in \eqref{eq:comparison-edit-energy}.

Take \(I_E\) to contain all rows within distance \(C'D_1\) of \(E\).
Choose the multiplier in \(d\) after the stencil constants, so that
the total stencil enlargement is less than \(d/2\).  It then reads no
marked part of an unselected dilated component: such marked parts
remain at distance at least \(d\) from \(E\).  On the other unchanged
sites the exterior estimate applies, and on replaced sites the sampled
chords are below \(t_0\).  Every row read by \(I_E\) is therefore
bounded by \(C_LHt_0^2=C_Lp_0^2\).  The count
\(|I_E|\le Cd^2n\) proves the second bound.
\end{proof}

\subsection{Integrating the covering identities}
\label{subsec:comparison-covering-algebra}

We now combine the row estimate and the replacement estimate.  This is
the step that avoids any pointwise division by a signed restricted sum.
First choose \(B_T\) and \(\vartheta\) so that
\begin{equation}
 (1-\vartheta)C_Ld^2<\frac{c}{4d^2},
 \qquad
 \frac{p_0^2}{d^2}\gg d^2B_T^2\log H+\log H+1.
 \label{eq:comparison-parameter-choices}
\end{equation}
These are the choices in \cite[Equations~(7.9)--(7.10)]{O3Continuum}.
The first inequality is achieved by Lemma~\ref{lem:comparison-row-moments};
the second follows by choosing \(P_0\) after the powers defining
\(d\) and \(B_T\), and then increasing \(H\).  We allow a fixed decrease
of \(c\) to reserve some direct energy for the original real-coordinate
integrations.

\begin{lemma}[Integrated gain from a specified edit]
\label{lem:comparison-edit-gain}
Fix the exact site sets of a selected family of original regions, let their union
be \(E\), and let \(n=|E|\).  Let \(\mathcal Q\) be any measurable set
of original flat configurations whose region construction gives those sets,
with any further inventory restrictions imposed before sampling.
Average the fillings of
Lemma~\ref{lem:comparison-flat-edits}, keeping the original geometry fixed,
and denote the filled configuration by \(z'\).  Then, for either time seam,
\begin{equation}
 \frac1{Z_{F,+}^{\rm end}}\int_{\mathcal Q}e^{X_F^\eta(z)}
       \mathbb E_{\rm fill\mid z}[\Xi_F^\eta(z')]\,d\nu(z)
 \le \exp(-c'p_0^2n/d^2).
\label{eq:comparison-edit-gain}
\end{equation}
This is the analogue of \cite[Equation~(7.15)]{O3Continuum}.
\end{lemma}

\begin{proof}
Only rows in \(I_E\) change.  Equations
\eqref{eq:comparison-remainder} and \eqref{eq:comparison-edit-energy}
give the exponent comparison
\[
 X_F^\eta(z)-X_F^\eta(z')
 \le C_Ln-cp_0^2n/d^2+\sum_{e\in I_E}Y_e(z').
\]
Retain a fixed fraction of the original direct energy on edges incident
on \(E\), decreasing \(c\).  The first inequality in
\eqref{eq:comparison-parameter-choices} then permits replacement of
the last coefficient by \(\vartheta\), at a cost of at most a fixed
fraction of the negative term.

Apply the uniform bound on the conditional fill density, and integrate
the original variables in \(E\) against the reserved direct weight.
When \(E\ne\Lambda\), choose a spanning forest anchored in the unchanged
exterior.  Each tree edge integrates a real increment, while each
remaining incident edge has an exponentially convergent integer sum.
The spacing \(P\) is bounded below, so these costs are at most
\(e^{C_Ln\log H}\), independently of \(P\).  For \(E=\Lambda\), one
tree root remains free and the cost is
\(Pe^{C_Ln\log H}\).  It cancels the factor \(1/P\) in the whole-torus
fill density.  Unchanged exterior data include their edge integers;
only integers incident on \(E\) are summed in this operation.

The final factor \(\Xi_F^\eta(z')\) is a complete positive flat-reference
sum.  We may therefore discard the restrictions defining \(\mathcal Q\)
and integrate the final configuration against the complete positive
reference density, retaining its flatness constraint.  The resulting
upper bound is
\[
 \exp\{C_Ln\log H-c''p_0^2n/d^2\}
 \frac1{Z_{F,+}^{\rm end}}
 \int_{\mathcal F}e^{X_F^\eta}\Xi_F^\eta
       e^{\vartheta\sum_{I_E}Y_e}\,d\nu.
\]
There is no injective-change-of-variables assertion in this argument:
the original and filled variables are integrated separately using the
conditional-density bound.  Apply \eqref{eq:R16},
\(|I_E|\le Cd^2n\), and the second inequality in
\eqref{eq:comparison-parameter-choices} to obtain
\eqref{eq:comparison-edit-gain}.
\end{proof}

\begin{proposition}[Integrated comparison on flat rectangles]
\label{prop:comparison-flat-integrated}
Fix a constant multiplier of the retained covering cap.  For sufficiently
large \(H\), let \(e^{X_F^\eta}\Xi_F^\eta\) be the terminal flat-reference
density constructed above, on any admitted dyadic rectangle whose
periods exceed the stated \(H\)-dependent constant.  Let
\(\widetilde\Xi^\eta\) be another mandatory covering sum with the same
inventory conventions and complete supports, allowing complex weights.
Use a common label list, setting absent weights equal to zero, and define
\[
 d_\ell=\|\widetilde\kappa_\ell-\kappa_{\ell,F}\|_\infty,
 \qquad
 b_\ell=\|\widetilde\kappa_\ell\|_\infty+
              \|\kappa_{\ell,F}\|_\infty.
\]
Suppose
\[
 \sup_x\sum_{\ell:x\in P_\ell}e^{64s_\ell}b_\ell
 \le C\exp(-p_0^{1/4}),
\]
with the fixed multiplier \(C\), and with the support and load bounds of
the retained class.  Then, for either \(\eta\),
\begin{equation}
 \frac{\displaystyle\int_{\mathcal F}e^{X_F^\eta}
              |\widetilde\Xi^\eta-\Xi_F^\eta|\,d\nu}
      {Z_{F,+}^{\rm end}}
 \le
 \exp\left(\sum_\ell d_\ell e^{2s_\ell}\right)-1.
 \label{eq:comparison-flat-cover}
\end{equation}
The constants are independent of the cutoff depth, the periods, their
aspect ratio, and \(P\).  The estimate also bounds the absolute difference
of the two integrals after insertion of any common measurable multiplier
of modulus at most one.
\end{proposition}

\begin{proof}[Proof of Proposition~\ref{prop:comparison-flat-integrated}]
We use the identities and forest estimate of
\cite[Sections~7.4--7.5]{O3Continuum}.  The relevant subsections are
``Exact covering identities and integration'' and ``Summing the forests.''

Fix the original regions and their inventories, and sample their fillings
once.  If \(I\) is a set of regions left unfilled, write \(z(I)\) for the
resulting configuration.  Then
\[
 D(z(I))=\bigsqcup_{C\in I}D_C.
\]
Write \(\Xi(I)=\Xi_F^\eta(z(I))\) and
\(\widetilde\Xi(I)=\widetilde\Xi^\eta(z(I))\).
No filled configuration is reclustered.

In a complete covering of these active inventories, join labels whenever
their supports meet the same active region.  Each connected group is a
macrolabel: its enlarged support consists of its label supports and all
active regions they meet, and its weight is the product of their weights.
It covers in full every active region it touches.  Distinct macrolabels
have disjoint enlarged supports.  For a site set \(S\) disjoint from the
active regions, let \(\Xi(I;S)\) be the complete-inventory reference sum
restricted to macrolabels avoiding \(S\).  This restricted sum may be signed.
Write \(w_U(z(I))\) for a macrolabel's product of reference weights, and
\(\widetilde w_U(z(I))\) for the corresponding product of comparison
weights; their enlarged supports are denoted by \(S_U\).

Selecting a compatible macrolabel family \(\mathcal A\), with combined
support \(S_{\mathcal A}\) and consumed regions \(B_{\mathcal A}\), leaves
the background sum
\(\Xi(I\setminus B_{\mathcal A};S_{\mathcal A})\).
Each selected macrolabel consumes every active region its support meets,
so \(S_{\mathcal A}\) is disjoint from the remaining active regions, as
required in the definition of this restricted sum.  Every remaining
label avoids all changed sites and incident integer data, so its value
and all eligibility tests are unchanged.  Its inventory loses exactly
the consumed inventories.  Expanding the avoidance indicators by
inclusion--exclusion, and expanding the difference of the two full
products, gives
\begin{align*}
 \Xi(I;S)
 &=\sum_{\substack{\mathcal A\ \mathrm{compatible}\\
                    S_U\cap S\ne\varnothing\ (U\in\mathcal A)}}
       (-1)^{|\mathcal A|}
       \prod_{U\in\mathcal A}w_U(z(I))\,
       \Xi(I\setminus B_{\mathcal A};S_{\mathcal A}),\\
 \widetilde\Xi(I)-\Xi(I)
 &=\sum_{\substack{\mathcal A\ \mathrm{compatible}\\
                    \mathcal A\ne\varnothing}}
       \prod_{U\in\mathcal A}
            [\widetilde w_U(z(I))-w_U(z(I))]\,
       \Xi(I\setminus B_{\mathcal A};S_{\mathcal A}).
\end{align*}
All macrolabels in these sums satisfy the full-inventory conditions just
described.  These are the cited exact identities.  Iterate the first in the
background factors arising from the latter.  Its empty-family term is a
complete reference sum, while every nonempty generation consumes an
active region.  Thus the recursion terminates at complete reference sums
\(\Xi_F^\eta(z')\).

Take absolute values only after these identities.  Telescoping a
macrolabel product difference marks one of its original labels, at cost
\(d_\ell\), while all other label factors cost \(b_\ell\).  The deterministic
forest encoding from the cited proof has original labels and exact
region site sets as vertices: internal edges describe each macrolabel,
and a second edge type joins a later macrolabel to an intersecting one
in the preceding generation.  Root each component at its marked label.
No region is repeated, because it is consumed at first use.  No original
label is repeated, because its fixed nonempty inventory becomes
ineligible once consumed.  The forest retains the generations, the
marked choices, and every evaluation configuration.  Thus the same
counting applies with the present incident-integer convention.
The common label list includes labels whose evaluated weight is zero;
all spin-dependent eligibility tests stay in the evaluated weights.
For a fixed encoded term the set \(\mathcal Q\) in
Lemma~\ref{lem:comparison-edit-gain} is consequently defined from the
original geometry and inventories before the fillings are sampled.

Lemma~\ref{lem:comparison-edit-gain} assigns a region of \(n\) sites
weight \(e^{-c'p_0^2n/d^2}\).  On the rectangular torus the number of
connected site animals of size \(n\) through a given site is at most
\(C^n\), with a constant independent of either period.  Let
\(R=C_LM_0^2\) be a fixed support-to-load bound, so that
\(|P_\ell|\le Rs_\ell\).  For \(H\) large,
\[
 \sup_x\sum_{E\ni x}e^{-c'p_0^2|E|/d^2}e^{2|E|}
 <\frac1{16R},\qquad
 \sup_x\sum_{\ell:x\in P_\ell}b_\ell e^{2s_\ell}
 <\frac1{16R}.
\]
Summing over an intersection site and the two edge types bounds the
one-child sum at a vertex of size \(s\) by \(s/4\), including an
exponential allowance for descendants.  Induction on tree depth bounds
all descendants by \(e^s\).  A marked root consequently costs at most
\(d_\ell e^{2s_\ell}\).  Dropping disjointness between the resulting rooted
trees and summing unordered nonempty root families gives
\[
 \sum_{j\ge1}\frac1{j!}
       \left(\sum_\ell d_\ell e^{2s_\ell}\right)^j,
\]
which is \eqref{eq:comparison-flat-cover}.  The support-hit bounds justify
passing from finite to countable label lists.  Positivity has been used
only for complete flat-reference sums at the filled configurations,
never for a restricted signed sum.  A common multiplier of modulus at
most one can be removed in this absolute estimate, proving the last
assertion as well.
\end{proof}

\begin{proof}[Completion of Proposition~\ref{prop:rg-comparison}]
Apply Proposition~\ref{prop:comparison-flat-integrated} with the Brownian
covering sum on \(\mathcal F\).  By \eqref{eq:cmp-terminal-distances}, the sum in \eqref{eq:comparison-flat-cover} is at most
\(C_HV e^{-c_H\beta}\), and the scalar-stripped exponents have the same
supremum discrepancy.  The direct, unbounded rows agree exactly at the
matched terminal couplings; all remaining exponent differences are
bounded.  The integrated estimate therefore also permits replacing the
Brownian exponent by the reference exponent.  Since \(V\) is polynomial
in \(\beta\), this gives an error \(O(e^{-c\beta})Z_{F,+}^{\rm end}\),
after decreasing \(c\).  Add \eqref{eq:comparison-nonflat}.

Finally insert the Brownian terminal phase in both flat integrals.  Its
modulus is one, with no bound required on its coefficient.  On the
reference side exact transport of winding through the observations
identifies it with the initial multiplier
\(\exp(i\phi\sum d_tv)\), for the corresponding real \(\phi\).
Reassemble the exact observations and remove each regulator's accumulated
bulk scalar.  The two scalar coefficients need not agree, but each is
independent of the periods and the time seam.  The same coefficient for
each regulator is therefore used for both seams and for time lengths
\(m_t\) and \(2m_t\).  The resulting scalar-stripped partition functions
satisfy \eqref{eq:R1}, proving Proposition~\ref{prop:rg-comparison}.
\end{proof}
\subsection{Pressure under a small fugacity change}
\label{subsec:comparison-pressure}

The Brownian spectral calculation requires control of the particle numbers
that maximize the transfer norm.  For this purpose we need an upper bound
on the change in the untwisted pressure when the logarithm of the fugacity
is shifted by either sign of $1/b$.  The geometry and the horizontal heat
time must remain fixed during this change.

\begin{lemma}[Pressure comparison at nearby fugacities]
\label{lem:comparison-nearby-fugacity}
Fix the blocking factor $L$ and the admission constants as above.  Let
$(b,r)$ be the matched Brownian parameters supplied by
\eqref{eq:R15}, with $b=\beta+O_L(1)$ and $r=1+O_L(1/b)$.
Consider the periodic rectangle with initial side lengths
\[
  X=L^N m_x,\qquad n=L^N m_t,\qquad
  m_x\asymp\beta(\log\beta)^2,\qquad
  m_t\asymp\sqrt\beta(\log\beta)^4,
\]
where the terminal lengths are dyadic and satisfy the admission conditions.
Write $V=m_xm_t$ and $V_{\mathrm{site}}=Xn=L^{2N}V$.
For $\sigma\in\{-1,1\}$ set
\[
  b'=b\exp\!\left(\frac{\sigma}{2b}\right),\qquad
  r'=r\frac{b'}b.
\]
Thus the fugacity $\lambda_{\mathrm{fug}}=b^2/2$ changes by
$\log(\lambda_{\mathrm{fug}}'/\lambda_{\mathrm{fug}})=\sigma/b$, while $r'/b'=r/b$.
In particular, the horizontal heat time per initial interval, the width
$W=Xr/b$, and the rectangle volume are unchanged.  For all sufficiently
large $\beta$, the actual untwisted Brownian partition functions satisfy
\begin{equation}
  \log Z^B(b',r')-\log Z^B(b,r)
  \le V_{\mathrm{site}}
       \left[b'r'-br+O_L\!\left(\frac1b\right)\right].
  \label{eq:R17}
\end{equation}
The error is uniform in $\sigma$ and in the indicated rectangle lengths,
with fixed constants in the two $\asymp$ conditions;
the assertion also holds with $m_t$ replaced by $2m_t$.
\end{lemma}

\begin{proof}
We first obtain the lower bound at the matched parameters, including the
phase in \eqref{eq:R1}.  In the initial variables, the flat reference is
the product of a Wilson integral and a Gaussian lift, summed over their
compatible windings.  If $(k_x,k_t)\in\mathbb Z^2$ is the winding pair,
the harmonic part of the Gaussian energy is
\[
  \frac{\pi^2}{2}
  \left(c_x\frac{m_t}{m_x}k_x^2+
        c_t\frac{m_x}{m_t}k_t^2\right),
  \qquad c_x,c_t=\beta+O_L(1).
\]
The Gaussian fluctuation integral is the same in every winding sector.
The corresponding Wilson integral has the center twists prescribed by
$(k_x,k_t)$, and is at most its untwisted value.  This last bound follows
by expanding the Wilson bonds in powers: after spin integration the
coefficients with no seam signs are nonnegative, and adding a center
seam only inserts signs.

Let $A_0$ denote the positive zero-winding contribution to the untwisted
reference partition function.  Since the phase has modulus one and is
identically one at zero winding, the total absolute contribution of all
other windings is at most $\varepsilon_\beta A_0$, where
\[
  \varepsilon_\beta
  \le \sum_{(k_x,k_t)\ne(0,0)}
       \exp\!\left\{-c\beta\left(
        \frac{m_t}{m_x}k_x^2+
        \frac{m_x}{m_t}k_t^2\right)\right\}=o(1).
\]
Indeed,
$\beta\min(m_t/m_x,m_x/m_t)\to\infty$ for the chosen rectangles.
Pairing opposite windings makes the phase-weighted reference integral
real.  It is therefore at least $(1-\varepsilon_\beta)A_0$, whereas the
reference integral with phase zero is at most
$(1+\varepsilon_\beta)A_0$.  These comparisons are unchanged when the
common reference volume scalar is removed.
In particular,
\[
 Z^{F(\phi)}_+\ge
 \frac{1-\varepsilon_\beta}{1+\varepsilon_\beta}Z^{F(0)}_+.
\]

Consequently the phase-weighted reference retains, up to a fixed factor,
the aligned-cap lower bound established above.  The integrated comparison
\eqref{eq:R1}, obtained from \eqref{eq:R15}, transfers this lower bound
to the matched Brownian integral after its volume scalars are removed.
At the terminal layer that integral is thus at least
\[
  cP\exp(-C_L V\log H),
\]
where $P=\pi R_0$ is the terminal lift period.

For the upper bound at $(b',r')$ we need only a short initial part of the
renormalization.  Initialize both runs with the same fixed $R_N$, and
perform $s=\lceil\log b/\log L\rceil=O_L(\log b)$ steps; in particular,
$s\ll N$.  At forward step $h\le s$ the depth index
is $j=N-h$ and $H_j\asymp b$.  The initial displacement in the displayed
couplings is bounded, and each coefficient step, with the prescribed
coordinate rescalings, costs at most $C_L$.  The displacement and the
accumulated corrections are therefore at most $C_L(1+h)$.  The admission
bands have room for this initial segment at both parameter values.

After $s$ steps the layer contains $L^{-2s}V_{\mathrm{site}}$ sites.
The spanning-tree integration bound and the class caps, with the phase
replaced by its absolute value, bound the remaining integral by
\[
  P_{N-s}\exp\!\left(
       C_L L^{-2s}V_{\mathrm{site}}\log b\right).
\]
Our choice of $s$ gives $L^{-2s}\log b\le C_L/b$.
The logarithm of this integral then costs
$O_L(1/b)$ per initial site.  The single factor $P_{N-s}$ has the same
property: $|\log P_j|=O_L(\log b)$ throughout the trajectory, and
$V_{\mathrm{site}}=L^{2N}V$ grows exponentially in $\beta$.

It remains to compare the extracted scalars on this short segment.
Let $S_s(b,r)$ denote their sum per initial site, including initialization
and the first $s$ steps.  A scalar extracted per output site at forward
step $h$ has weight $L^{-2(h+1)}$ in this sum.  By \eqref{eq:R14}, the
initial scalar difference is
\[
  b'r'-br+O_L(1/b).
\]
Here the changes in $\log g$ and in the bounded Lipschitz initialization
term are $O_L(1/b)$, and the two scalar remainders have that size as well.
On the subsequent short segment, the bounds $H_j\asymp b$ and
$C_L(1+h)$ on coefficient drift show that the shape, $\log g$, and
$\chi$ arguments in \eqref{eq:R11} differ by at most
$C_L(1+h)/b$.  The rescaling schedule and the type of free history are
the same in the two Brownian runs.  The Lipschitz bound for the displayed
terms in \eqref{eq:R11}, together with the absolute $O_L(1/H_j)$ bounds
on its two remainders, consequently gives
\begin{align*}
  S_s(b',r')-S_s(b,r)
  &=b'r'-br+
     O_L\!\left(\frac1b\right)
       \left(1+\sum_{h=0}^{s-1}(1+h)L^{-2(h+1)}\right)\\
  &=b'r'-br+O_L(1/b).
\end{align*}

For the matched run, the scalars extracted after these $s$ steps also
have total absolute size $O_L(1/b)$ per initial site.  Indeed
\eqref{eq:R11} bounds each by $C_L\log b$ per output site, and their
remaining volume weights sum to $O_L(L^{-2s})$.
The terminal lower bound above has logarithm bounded below by
$\log(cP)-C_LV\log H$, whose negative part, divided by
$V_{\mathrm{site}}$, is again $O_L(1/b)$.  We have therefore proved
\[
  \frac{\log Z^B(b',r')}{V_{\mathrm{site}}}
       \le S_s(b',r')+O_L(1/b),\qquad
  \frac{\log Z^B(b,r)}{V_{\mathrm{site}}}
       \ge S_s(b,r)-O_L(1/b).
\]
Subtracting and using the scalar comparison proves \eqref{eq:R17}.
The argument uses the full matched trajectory only for the lower bound;
the displaced parameters require just the admitted short segment.
\end{proof}
